\section{Discussion}
\label{sec:discussion}

\subsection{Why SFT on 7B Q8 worked where DPO/ORPO on 1.5B 4-bit did not}
\label{sec:why-it-worked}

Before arriving at the Phase~1 recipe reported here, this project ran
preference-optimization experiments on Qwen~2.5\,1.5B at 4-bit
quantization. Both DPO \cite{rafailov2023dpo} and ORPO \cite{orpo}
produced apparently strong validation signals -- DPO converged to
1.0 accuracy with a 7.5-nat margin between chosen and rejected
responses -- while generations from the trained adapters were nearly
byte-identical to the base for most prompts. For ORPO, generations for
three of four spot-checked prompts were \emph{literally} byte-identical
to base.

This is the well-documented asymmetric-reward failure mode of DPO: the
likelihood gap can be widened by pushing the rejected response down rather
than by pulling the chosen response up, in which case the policy
moves very little in the regions of the output distribution that
inference actually samples from. Three factors plausibly compounded the
problem in our 1.5B 4-bit setting:

\begin{enumerate}
  \item \textbf{Capacity.} A 1.5B model has less unused representational
        room to write a new register on top of the existing
        helpful-assistant policy. Stylistic transfer competes more
        directly with the model's other behaviors.
  \item \textbf{Quantization.} 4-bit weight quantization, while necessary
        for memory at small budgets, limits the precision with which the
        adapter can express small but consequential shifts in
        next-token distribution. We hypothesize, without proving, that
        8-bit is qualitatively friendlier to LoRA-style adaptation.
  \item \textbf{Objective.} Preference objectives are a strange first
        teaching signal when the base model has effectively zero
        probability mass on the target register. SFT directly maximizes
        the likelihood of target-register sequences and so injects mass
        where preference optimization can only reweight it.
\end{enumerate}

The Phase~1 recipe deliberately addressed all three: a larger base
(7B), a gentler quantization (Q8), and SFT as the first training step.
We do not claim that DPO is generally inferior to SFT for register
transfer; we claim that for our setting -- small handcrafted corpus,
small open-weights base, on-device training -- starting with SFT was
the right call, and a useful negative result for practitioners
considering preference optimization with similar constraints.

\subsection{The Augustinian gap}
\label{sec:augustinian-gap}

The most striking dimension-level failure in Table~\ref{tab:rubric} is
\emph{Augustinian voice}, which moves from 0 to only 3 out of 30. The
cause is a corpus imbalance: of the 342 source chunks, the Summa
dominates by chunk count, and the Claude teacher's pair generator was
more confident producing Summa-form outputs than Augustinian-form ones.
The 83-pair training mix is consequently heavily Summa-weighted, and
the model has not been shown enough Augustinian register to internalize
its surface markers (second-person address to God, autobiographical
framing, scriptural allusion).

The Phase~2 plan addresses this by oversampling Confessions and City of
God chunks during pair generation and by adding an explicit
Augustinian-voice instruction variant in the teacher's system prompt.

\subsection{Ethics}
\label{sec:ethics}

The trained model is not a theological authority and must not be
treated as one. Two concrete risks deserve calling out:

\begin{itemize}
  \item \textbf{Hallucinated CCC citations.} The fine-tuned model
        confidently emits ``CCC~$n$'' references with the surface form of
        ground truth. The rubric counts these as \emph{grounding}; it
        does not verify them. Spot checks show that a non-trivial fraction
        of paragraph numbers do not correspond to the actual content of
        the cited paragraph. Users must verify against the actual
        Catechism before trusting any specific citation.
  \item \textbf{Doctrinal authority.} The model speaks in a voice
        culturally associated with magisterial authority. It is not
        magisterially authoritative, has no episcopal review, and can
        confidently err. The repository's \texttt{README.md} carries a
        prominent NOTICE to this effect; any publication should
        reproduce that disclaimer.
\end{itemize}

\subsection{Limitations}
\label{sec:limitations}

\paragraph{Evaluation set size.} Ten held-out prompts is a small sample.
Variance estimates on per-dimension and total scores are therefore
generous; the +49-point delta is large relative to noise but we have
not bootstrapped a confidence interval.

\paragraph{Rubric coarseness.} A regex-based rubric measures lexical and
structural surface form. It cannot detect theological correctness,
internal consistency, or the difference between a well-formed objection
that resolves coherently and one whose ``I answer that'' contradicts the
``Reply to Objection 1.'' The rubric is a fast iteration tool, not a
judgment.

\paragraph{No human evaluation.} We did not run a human rater study. A
small expert eval (e.g., three theologically literate raters across 30
prompts) would tighten claims about scholastic register quality and
catch problems the rubric cannot.

\paragraph{Adversarial prompts.} We did not test the model on
intentionally adversarial prompts (jailbreaks, theological provocations,
out-of-domain queries that invite hallucinated citations). On-domain
behavior is the focus here; full safety evaluation is future work.

\paragraph{Fair-use posture on CCC.} Catechism text is used in source
material, in teacher-generated answer pairs, and in adapter weights
that arguably constitute a derivative work. We treat the LoRA weights
themselves as non-redistributable by default; only code, recipe, and
aggregate metrics are released. A future public model release would
need a careful review and would likely have to restrict CCC content to
ideas rather than verbatim text.
